% ==============================================================================
% AstroLib Academic Typography: Formula Stress Test Suite
% Covers AMS math, Greek letters, Blackboard bold, Fraktur, Calligraphic,
% multi-integrals, matrices, piecewise cases, sub/superscript optical alignment.
% ==============================================================================

\subsection*{1. 数学字符集与字形覆盖测试}

% Blackboard Bold & Calligraphic & Fraktur
实数域 $\mathbb{R}$、复数域 $\mathbb{C}$、整数环 $\mathbb{Z}$、有理数域 $\mathbb{Q}$ 与自然数集 $\mathbb{N}$。
希尔伯特空间 $\mathcal{H}$、傅里叶变换算子 $\mathcal{F}$、拓扑流形族 $\mathcal{M}$。
李代数 $\mathfrak{g}$、酉群李代数 $\mathfrak{su}(2)$ 与典型理想 $\mathfrak{a} \subset \mathfrak{g}$。

% Greek Alphabet (Uppercase & Lowercase)
小写希腊字母: $\alpha, \beta, \gamma, \delta, \epsilon, \varepsilon, \zeta, \eta, \theta, \vartheta, \iota, \kappa, \lambda, \mu, \nu, \xi, \pi, \varpi, \rho, \varrho, \sigma, \varsigma, \tau, \upsilon, \phi, \varphi, \chi, \psi, \omega$。\\
大写希腊字母: $\Gamma, \Delta, \Theta, \Lambda, \Xi, \Pi, \Sigma, \Upsilon, \Phi, \Psi, \Omega$。

\subsection*{2. 复杂极限、微积分与多重积分公式}

\begin{equation}
  \lim_{n \to \infty} \sum_{k=1}^n \frac{n}{n^2 + k^2} = \int_0^1 \frac{1}{1 + x^2} \,\mathrm{d}x = \left. \arctan x \right|_0^1 = \frac{\pi}{4}.
\end{equation}

高斯散度定理与麦克斯韦电磁场旋度算子：
\begin{equation}
  \oiint_{\partial \Omega} \mathbf{F} \cdot \mathrm{d}\mathbf{S} = \iiint_{\Omega} (\nabla \cdot \mathbf{F}) \,\mathrm{d}V, \quad \nabla \times \mathbf{E} = -\frac{\partial \mathbf{B}}{\partial t}.
\end{equation}

\subsection*{3. 大括号分块矩阵与分段函数}

\begin{equation}
  \mathbf{M}_{\text{block}} = \left(
    \begin{array}{cc|cc}
      a_{11} & a_{12} & 0 & 0 \\
      a_{21} & a_{22} & 0 & 0 \\
      \hline
      0 & 0 & \lambda_1 & 1 \\
      0 & 0 & 0 & \lambda_1
    \end{array}
  \right), \quad
  f(x, y) = \begin{cases}
    \frac{x^3 y - x y^3}{x^2 + y^2}, & (x, y) \ne (0, 0) \\
    0, & (x, y) = (0, 0)
  \end{cases}
\end{equation}

\subsection*{4. 光学上下标与微观韵律基线}

连续行内上下标与算子对比：
设 $T^{\mu\nu}_{\rho\sigma}$ 为四阶能动张量，其联络导数为 $\nabla_\lambda T^{\mu\nu}_{\rho\sigma} = \partial_\lambda T^{\mu\nu}_{\rho\sigma} + \Gamma^\mu_{\lambda\alpha} T^{\alpha\nu}_{\rho\sigma} - \Gamma^\alpha_{\lambda\rho} T^{\mu\nu}_{\alpha\sigma}$。
此时在中文连续排版行内，上下标不应与上一行底线或下一行顶线发生任何字形碰撞或挤压。
